% Generated draft: requires mathematical and source review.
\documentclass[11pt]{article}
\usepackage[margin=25mm]{geometry}
\usepackage{fontspec}
\setmainfont{Latin Modern Roman}
\usepackage{amsmath,amssymb,mathtools,unicode-math}
\setmathfont{Latin Modern Math}
\usepackage{enumitem,needspace,xurl,hyperref}
\setlength{\parindent}{0pt}
\setlength{\parskip}{5pt}
\setlength{\emergencystretch}{3em}
\newcommand{\N}{\mathbb{N}}
\newcommand{\Z}{\mathbb{Z}}
\newcommand{\Q}{\mathbb{Q}}
\newcommand{\R}{\mathbb{R}}
\newcommand{\C}{\mathbb{C}}
\newcommand{\F}{\mathbb{F}}
\newcommand{\A}{\mathbb{A}}
\newcommand{\PP}{\mathbb{P}}
\newcommand{\E}{\mathbb{E}}
\newcommand{\eps}{\varepsilon}
\newcommand{\dd}{\,\mathrm{d}}
\begin{document}
\section*{UnsolvedMath ID: 11000132}
\textbf{Source problem code:} catalogue entry
\subsection{Definitions and mathematical statement}
Fix the surface $F$ and the decomposition, reference arcs, and twist conventions used for Dehn--Thurston coordinates in the source. For an admissible integral coordinate tuple
\[
 u=((m_1,t_1),\ldots,(m_N,t_N)),
\]
let $\alpha(u)$ be the corresponding weighted family of pairwise disjoint curves and arcs. In the usual interior-pants-curve convention, $m_i$ is an intersection measure and $t_i$ a twist coordinate; boundary/arcs and zero-intersection cases require the source's additional conventions and admissibility conditions. Not every arbitrary integer tuple represents such a family.

\textbf{Problem 2.} Give a tractable, more nearly closed-form expression in these or other coordinates for the number of components of $\alpha(u)$. The source explicitly distinguishes this request from the known serial train-track-splitting algorithm that recovers a multiply weighted curve from an integral train-track measure.

The component-count convention must be fixed. If
\[
 \alpha(u)=w_1\alpha_1+\cdots+w_q\alpha_q
\]
is written with distinct nonisotopic support curves or arcs and positive integral weights, counting support components gives $q$, while interpreting weights as parallel multiplicities gives $\sum_{j=1}^q w_j$. These are different coordinate functions; the excerpt does not state which is meant or how parallel arcs are grouped. The target is an expression for the intended function $C(u)$ after that convention is recovered, not an invented replacement formula.

The source additionally asks for a description in Dehn--Thurston coordinates on $\operatorname{MF}(F)$ of the loci denoted $\mathcal A'(F)$ or $\operatorname{Arc}'(F)$. Here $\operatorname{MF}(F)$ is the relevant measured-foliation space, but the precise foliation convention and both primed loci are not defined in this record. It suggests quantum path-ordering techniques [5,6] or fermionic statistical physics [57,29] as possible approaches and refers to Problem 3. These are speculative methods and an unreproduced cross-reference, not a solution or a supplied computational complexity guarantee.
\subsection{Short English statement}
Find a tractable coordinate formula for the component count of an integral weighted curve-and-arc family, beyond reconstructing it by repeated train-track splitting. Clarify admissible coordinates and whether multiplicities or only distinct support components are counted.
\subsection{Sources}
\begin{itemize}
\item Local UnsolvedMath registry, permanent ID 11000132, source code catalogue entry; no proposer or proposal year supplied.
\item Farb, edited volume \emph{Problems on Mapping Class Groups and Related Topics}, chapter context ``Probing mapping class groups using arcs,'' Problem 2, PDF page 113; \url{https://www.math.uchicago.edu/~farb/papers/mcgbook.pdf}. The registry records PDF-font-structure extraction and access on 2026-07-29; the inventory distinguishes 329 labeled headings by source page.
\item Source references [6] for splitting and [5,6,57,29] for suggested quantum/statistical techniques, without full bibliography entries. The research summary cites Penner--Harer, \emph{Combinatorics of Train Tracks} (1992), as related background rather than as a supplied closed-form answer.
\end{itemize}
\end{document}
